% This file is part of Astatistics.
% Copyright (C) 2011
% 	Eurogaran Informatica, S.L.
% 
% Astatistics is free software: you can redistribute it and/or modify
% it under the terms of the GNU General Public License as published by
% the Free Software Foundation, either version 3 of the License, or
% (at your option) any later version.
% 
% Astatistics is distributed in the hope that it will be useful,
% but WITHOUT ANY WARRANTY; without even the implied warranty of
% MERCHANTABILITY or FITNESS FOR A PARTICULAR PURPOSE.  See the
% GNU General Public License for more details.
% 
% You should have received a copy of the GNU General Public License
% along with Astatistics.  If not, see <http://www.gnu.org/licenses/>.
% 
% See the file COPYING for copying conditions.

\documentclass[spanish,12pt]{book}
\usepackage[utf8]{inputenc}
\usepackage{babel}
\usepackage{amssymb}
\usepackage{enumerate}
\usepackage{a4wide}
\setcounter{secnumdepth}{0}
\setcounter{tocdepth}{3}
\setlength\parskip{\medskipamount}
\setlength\parskip{\baselineskip}
\setlength\parindent{0pt}
\pagestyle{plain}
\usepackage[svgnames]{xcolor}
\usepackage{graphicx}
\usepackage{cjhebrew}
\usepackage{pdfpages}
\usepackage{listings}

\begin{document}

%\title{\center{\includegraphics{artwork/logo_blc.ps}}
%	\vspace{1cm}
%	\linebreak Manual de usuario}
\title{Manual de usuario}
\author{\includegraphics{artwork/logo-mini-oval-peq.ps}}
\maketitle
%\includepdf[pages={1}]{tapa.pdf}
%\thispagestyle{empty}
\cleardoublepage

\pagenumbering{roman} \tableofcontents \cleardoublepage \newpage \pagenumbering{arabic}

\chapter{Introducción}
Astatistics es un avanzado sistema de informes y estadísticas para Asterisk accesible desde el navegador web diseñado para ser completamente flexible y adaptable.

A diferencia de otros programas similares, Astatistics permite al usuario definir qué estadísticas son necesarias en cada instalación y cómo presentarlas al usuario final.

Dispone de un potente sistema de usuarios, permisos y roles, un potente motor de diseño de menús jerárquicos, y un óptimo motor de acceso a los datos en tiempo real, así como de funciones de exportación de datos a diferentes formatos.

Está basado en el framework de programación Catalyst y es muy fácil de instalar y configurar.

Este manual, cuyo fin es facilitar el uso del programa, está destinado tanto al adminsitrador del sistema, encargado de la instalación y mantenimiento de Astatistics, como a tres perfiles de usuario distintos dentro de Astatistics: el administrador, el diseñador y el usuario.

El administrador es el perfil de usuario que mantiene la plataforma y tiene permisos totales en Astatistics. Puede diseñar y visualizar plantillas tanto diseñadas por él como por otros usuarios, además de gestionar las propias cuentas de usuario.

El diseñador es el perfil de usuario dedicado al diseño de plantillas de estadísticas. Podrá diseñar o visualizar aquellas plantillas o estadísticas para las cuales tenga permisos de acceso.

El usuario es el perfil que realizará las consultas finales. Tan solo podrá visualizar resultados para aquellas plantillas que tenga acceso concedido mediante el sistema de permisos y roles.

Astatistics se distribuye bajo licencia GPL de GNU, por lo que se trata de software libre. Puede modificar y redistribuir este software bajo los términos de dicha licencia. Si desea conocer más detalles consulte el fichero COPYING de nuestra distribución de Astatistics o visite {\textless}http://www.gnu.org/licenses/\textgreater.

\chapter{Instalación de Astatistics (Administrador del sistema)}
En las intrucciones que se detallan a continuación el símbolo `\$' precede a instrucciones que no requieren privilegios de root, mientras que el símbolo `\#' precede a las que sí requieren dichos privilegios.

\section{Requisitos de instalación}
Para iniciar la instalación de astatistics se requiere los siguientes requisitos:
\begin{itemize}
\item Un sistema GNU/Linux con Asterisk completamente funcional del cuál se extraerá la información.
\item El módulo res\_mysql de asterisk. Se encuentra en el paquete de Asterisk Addons.
\item El servidor de bases de datos MySQL (probado con la versión 5.1).
\item El intérprete del lenguaje PERL y sus módulos base correctamente instalados y configurados. Este requisito lo cumplen la mayoría de las distribuciones GNU/Linux con una instalación estándar.
\item CPAN, un sistema de gestión de módulos de PERL.
\item CPANPLUS, otro sistema de gestión de módulos de PERL.
\item Una conexión a Internet completamente funcional para el acceso a CPAN.
\item El programa make de GNU (se instala automáticamente al usar el script de instalación para Debian GNU/Linux).
\item Un compilador correctamente configurado, como GCC, necesario para la compilación de algunos módulos de PERL (se instala automáticamente al usar el script de instalación para Debian GNU/Linux).
\item El programa wkhtmltopdf para exportación a PDF (se instala automáticamente al usar el script de instalación para Debian GNU/Linux).
\item Los paquetes libtest-nowarnings-perl libexpat-dev pkg-config libcairo-dev (nomenclatura en Debian GNU/Linux) o similares en el caso de otras distribuciones correctamente instalados  (se instalan automáticamente al usar el script de instalación para Debian GNU/Linux).
\end{itemize}

\section{Configurando MySQL y Asterisk}
Primero se ha de configurar el módulo res\_mysql.so editando el fichero de configuración res\_mysql.conf del directorio de configuración de Asterisk. En dicho fichero, en la sección ``general'' se deberá de introducir la información para acceder a la base de datos de Asterisk en el MySQL.

Se deberá de crear la base de datos de Asterisk y las tablas. Para ello se incluye un diseño en el fichero asterisk-mysql.sql del directorio ``db'' del paquete de instalación. Los pasos a realizar serían los siguientes (modificar según conveniencia):
\begin{enumerate}
\item Accedemos a mysql y hacemos:\\\\mysql\textgreater create database asterisk;\\Query OK, 1 row affected (0.01 sec)\\mysql\textgreater quit\\Bye\\\\Con esta operación habremos creado la base de datos ``asterisk''.
\item Cargamos el contenido del fichero asterisk-mysql.sql desde el directorio raíz del paquete:\\\\\$ mysql -u root -ppassword -e "$\backslash$. ./db/asterisk-mysql.sql"\\\\Con ello habremos creado las tablas dentro de la base de datos.
\item Definimos y concedemos permisos de acceso al usuario que hemos configurado anteriormente en el res\_mysql.conf:\\\\\$ mysql -u root -ppassword -e "grant all privileges on asterisk.* to 'usuario'@'localhost'
\end{enumerate}

Ahora ya tendremos creada la base de datos y configurado el acceso desde Asterisk.

A su vez es necesario configurar la salida del CDR de asterisk, así como del queue\_log y queue\_memebers a una base de datos MySQL de la siguiente forma:
\begin{enumerate}
\item Generamos un fichero (o modificamos el existente) ``/etc/asterisk/cdr\_mysql.conf'' con el siguiente contenido:\\\\``global''\\hostname=localhost\\dbname=asterisk\\password=passwd\\user=asterisk\\\\Modificando los datos con los que hemos usado anteriormente.
\item Modificamos los valores siguientes en la sección ``settings'' del fichero /etc/asterisk/extconfig.conf:\\\\queue\_members =\textgreater mysql,general,queue\_members\\queue\_log =\textgreater mysql,general,queue\_log
\end{enumerate}

Seguidamente haremos una recarga de la configuración de asterisk de la siguiente manera:\\\\\# asterisk -rx "reload"

Tras este proceso Asterisk debería de estar volcando ya la información de las llamadas y colas a las tablas creadas en MySQL. Se han incluído, a su vez, tablas de información de estado en tiempo real para la monitorización en vivo y otras facilidades.

\section{Instalando requisitos de módulos PERL}
Debido a un bug en la última versión de CPAN de DBIx::Class la instalación automática provista por Catalyst no es operativa en esta versión, por lo que habrá que hacer unos pasos preliminares. En el caso de estar instalando en Debian GNU/Linux Squeeze se puede realizar toda esta preparación tan solo ejecutando script/install-squeeze.sh con un usuario con permisos de ``sudo'' o bien como usuario root. En caso contrario:
\begin{itemize}
\item Instalamos la última versión de CPAN y CPANPLUS. Dependiendo de nuestro sistema PERL se deberían de utilizar diferentes instrucciones. Las más comunes para un sistema Debian GNU/Linux son:\\\\\# cpan CPAN\\\# cpan CPANPLUS\\\\Deberemos de responder a las preguntas que se nos hagan según nuestra conveniencia.
\item Instalamos el módulo Params::Validate:\\\\\# cpan Params::Validate
\item Instalamos Moose:\\\\\# cpan Moose
\item Instalamos el DBIx::Class con opción \-\-force para obviar tests que no funcionan en esta versión (la instalación suele ser completamente funcional, puesto que el error que se produce está en el test de acceso a SQLite, pero no en las funciones que testea):\\\\\# cpan -fi DBIx::Class
\item Finalmente podemos arrancar la instalación automática de Catalyst:\\\\\$ perl Makefile.PL build\\\# make
\end{itemize}

Tras estos pasos, y si todo ha ido bien, dispondremos de todas las dependencias de Astatistics instaladas correctamente en el equipo.

Para configurar el acceso a datos de Astatistics habrá que modificar la base de datos, el usuario y la contraseña del fichero script/create\_asterisk\_model.sh y luego ejecutarlo. Con ello el módulo DBIx::Class de PERL configurará todos los accesos a la base de datos Asterisk automáticamente.

Para probar el correcto funcionamiento de Astatistics podemos lanzar una instancia del servidor de desarrollo utilizando ``\$ script/astatistics\_server.pl''.

\section{Deployment}
El framework Catalyst puede ser ejecutado con el servidor incorporado, como en el ejemplo anterior, pero dicho servidor está diseñado para el desarrollo y tiene algunas deficiencias en producción, como que solo puede atender una petición simultáneamente.

Para un entorno en producción recomendamos utilizar el servidor Apache con FastCGI que tiene un comportamiento mejor y más rápido. A continuación detallamos los pasos a realizar para configurar Astatistics de esta forma:

Añadir un vhost con la siguiente configuración en Apache2 (consultar documentación):

\begin{lstlisting}
<VirtualHost 0.0.0.0:80>
	ServerAdmin webmaster@localhost
	
	ServerName astatistics
	DocumentRoot /path/Astatistics/root
	<Directory />
		Options FollowSymLinks
		AllowOverride None
	</Directory>
	<Directory /var/www/>
		Options Indexes FollowSymLinks MultiViews
		AllowOverride None
		Order allow,deny
		allow from all
	</Directory>

	Alias /tmp /tmp/wkhtmltopdf

	Alias /static /path/Astatistics/root/static
	FastCgiServer /path/Astatistics/script/astatistics{\_}fastcgi.pl -processes 3 -idle-timeout 120
	Alias /	/path/Astatistics/script/astatistics{\_}fastcgi.pl/

	ErrorLog /var/log/apache2/error.log

	\# Possible values include: debug, info, notice, warn, error, crit,
	\# alert, emerg.
	LogLevel warn

	CustomLog /var/log/apache2/access.log combined
</VirtualHost>
\end{lstlisting}

La palabra /path habrá de ser sustituida por la ruta donde se habrá ubicado el árbol principal de Astatistics, por ejemplo /usr/local/astatistics.

En el listado anterior se ha utilizado el nombre de vhost ``astatistics'', por lo que deberemos apuntar nuestro navegador hacia dicho nombre para accederlo. Podemos configurar nuestro DNS incluyendo el nombre escogido o, para acceder desde nuestro sistema local, añadir ``astatistics'' al fichero /etc/hosts con nuestra IP local o de loopback (127.0.0.1 en la mayoría de los casos).

Asegúrese de que el directorio que contiene la base de datos de astatistics (SQLite) y el fichero de dicha base de datos (astatistics.db) son accesibles y tienen permiso de escritura por el usuario utilizado por el servidor web (en Debian www-data). En el ejemplo siguiente se hace dando permiso de escritura para el grupo www-data:

\begin{lstlisting}
# chgrp www-data db/
# chgrp www-data db/astatistics.db
# chmod 775 db/
# chmod 664 db/astatistics.db
\end{lstlisting}

Para permitir que se ejecute correctamente la exportación a PDF, el paquete wkhtmltopdf ha de ejecturarse desde un servidor X. Para ello hemos añadido un wrapper denominado wkhtmltopdf-catalyst en el directorio script/ que debería de ser copiado a /usr/local/bin:

\begin{lstlisting}
# cp script/wkhtmltopdf-catalyst /usr/local/bin
\end{lstlisting}

Dicho script utiliza el binario xvfb-run, por lo que dicho ejecutable debería de estar instalado en el sistema. En debian esto se consigue instalando el paquete xvfb.

Tras reiniciar el servidor Apache con ``/etc/init.d/apache2 restart'' el sistema Astatistics debería de ser accesible correctamente, en nuestro ejemplo apuntando el navegador a la dirección ``http://astatistics''.

Para propagar la configuración de nuestra base de datos de Asterisk y hacerla utilizable por Astatistics se habrá de ajustar y luego ejecutar el script ``script/create\_asterisk\_model.sh''.

Según la instalación también podría ser necesario regenerar el modelo de la base de datos SQLite de Astatistics revisando y ejecutando el script ``script/create\_model.sh''.

Tras estos pasos ya se podrá acceder y probar el sistema Astatistics.

\chapter{Astatistics: Administración de usuarios}
Inicialmente el sistema viene preconfigurado con 3 usuarios: admin, designer y viewer, todos con contraseña 4321 por omisión.

Los usuarios de astatistics tienen diferentes tipos de permisos de sistema según los roles asociados a los mismos. Hay dos roles de sistema básicos:

\begin{itemize}
\item {\bf superuser}: Un usuario con este rol podrá acceder a todas las funciones del programa, incluyendo diseño de plantillas de cualquier usuario y administración de los otros usuarios. El usuario ``admin'' de la configuración por omisión tiene este rol activado.
\item {\bf can\_design}: Un usuario con este rol podrá diseñar nuevas plantillas y accederlas tras su creación, tanto para visualizar resultados como para modificar su diseño. El usuario ``designer'' de la configuración por omisión tiene este rol activado.
\end{itemize}

Además de estos roles de sistema, cualquier usuario con el rol superuser activado podrá crear nuevos roles y asignarlos a otros usuarios. De esta forma se pueden establecer los niveles de acceso que sean necesarios para cada usuario, pues se puede definir que roles son necesarios para acceder una plantilla.

La administración de usuarios se llevará a cabo desde la opción de menú ``Users'' y la administración de roles desde el menú ``Roles''.

Al acceder a dichos menús se mostrará una lista según el nivel de acceso del usuario. Un usuario común tan solo puede modificar su información (identificador y contraseña), mientras que un usuario con superuser activado podrá también modificar la información de otros usuarios, así como activar o desactivar roles para los mismos.

Se pueden definir tantos usuario o roles como sean necesarios.

\chapter{Astatistics: Diseño y gestión de plantillas y menús}
\section{Estructura de menús}
Los informes de Astatistics se ubican en un sistema gerárquico de menús para con objeto de poder organizar la información según las necesidades de cada instalación.

Los menús pueden tener el nombre que se desee y pueden ser cíclicos (un menú puede contenerse a sí mismo o a menús de los que cuelga. Así mismo los informes pueden aparecer en más de un menú.

Un menú, así como un informe, está sujeto al sistema de roles, por lo que solo será mostrado a aquellos usuarios que tengan permiso para visualizarlo.

\section{Las plantillas}
Una plantilla es un informe descrito con el lenguaje de plantillas de Astatistics, pudiendo solicitar parámetros al usuario antes de su visualización. Las plantillas sirven para definir qué tipo de informes son necesarios en cada instalación.

\section{Base de datos de Asterisk}
La fuente de la que se extrae la información para mostrar en los informes es la base de datos de Asterisk anteriormente creada. Se puede visualizar cualquier tipo de información sobre dicha base de datos, aunque normalmente se consultarán dos tablas:

\begin{itemize}
\item {\bf cdr}: Esta tabla almacena la información de las llamadas. Su estructura se define en el fichero ``/etc/asterisk/cdr\_custom.conf''.
\item {\bf queue\_log}: Esta tabla almacena los eventos que se suceden relacionados con la gestión de colas de llamadas (queues). Es una tabla especial puesto que tan solo almacena eventos. Para poder consultar estadísticas sobre dichos eventos se ha añadido una capa lógica en Astatistics que provee al diseñador de plantillas de la capacidad de consultar eventos lógicos como los tiempos de espera de los operadores, la duración de las llamadas, etc. Para acceder a los datos lógicos (virtuales) de dicha tabla se ha de especificar como tabla ``queuelog\_virt''.\end{itemize}

Además de estas dos tablas, en la definición de las tablas de Asterisk se añadieron varias tablas cuyos datos se llenan mediante triggers en las dos anteriores, facilitando la consulta de estados en tiempo real y otros datos adicionales:

\begin{itemize}
\item {\bf agent\_status}: Tabla que informa sobre el estado de los agentes en tiempo real.
\item {\bf call\_status}: Tabla que informa sobre el estado de las llamadas en curso.
\item {\bf queue\_log\_processed}: Tabla que contiene la misma información que la tabla queue\_log de una forma más ordenada para su consulta directa, así como con la marca de tiempo en formato de fecha de la base de datos, para poder hacer ordenaciones.
\end{itemize}

También se dispone de dos tablas ``queue\_members'' y ``queue\_table'' para el manejo en tiempo real de las colas de llamadas.

\section{Definición de plantillas y menús}
Una plantilla o menú tiene los siguientes datos para su definición:
\begin{itemize}
\item {\bf Class}: Group o Template, el primero para indicar que se trata de un menú (grupo de plantillas y menús) y el segundo para indicar que se trata de una plantilla.
\item {\bf Name}: Nombre o identificador, que es el texto que saldrá en la estructura de menús para identificar dicho menú o plantilla.
\item {\bf Roles}: Roles que permiten el acceso a dicha plantilla o menú. Si un usuario tiene permiso para acceder a cualquiera de los roles marcados podrá acceder a este ítem.
\end{itemize}

Las plantillas, además tienen los siguientes datos que no se usan en el caso de menús:
\begin{itemize}
\item {\bf Visual}: Tipo de visualización que se generará entre listado, gráfico de barras, gráfico de líneas, gráfico de áreas, gráfico circular (queso), gráfico de burbujas, gráfico de puntos y gráfico de área polar.
\item {\bf Conditions}: Dato donde se definen las condiciones para generar el informe, descrito en detalle más adelante. Consta de diferentes parámetros separados por punto y coma ``;''.
\item {\bf Format}: Parámetros relacionados con la representación final del informe, separados con punto y coma ``;''.
\end{itemize}

Los parámetros de Conditions y Format se especifican como una serie de asiganciones separadas por punto y coma ``;''. Un ejemplo sería el siguiente:
\begin{lstlisting}
table=agent_status;
show=agentId,agentName,agentStatus,timestamp,callid,queue;
refresh=#Refresco#5#;''
\end{lstlisting}

Los parámetros para Conditions son los siguientes:
\begin{itemize}
\item {\bf table}: Tabla que se usará para la consulta, tal como se denomina en la base de datos, exceptuando el caso de la tabla virtual queuelog\_virt.
\item {\bf show}: Campos de dicha tabla o funciones que se mostrarán en el informe. En el caso de gráficas se describen como coordenadas con tipo, campo y nombre, como: ``x:date:calldate:Fecha,y:COUNT:calldate:Llamadas''.\\\\Se pueden utilizar funciones agregadas de la base de datos sobre los campos de una tabla. Las funciones permitidas sobre tablas reales son las utilizadas en el lenguaje SQL (normalmente MySQL):
	\begin{itemize}
	\item MAX(): Valor máximo.
	\item MIN(): Valor mínimo.
	\item AVG(): Media aritmética.
	\item SUM(): Total de la suma de valores.
	\item COUNT(): Número de resultados.
	\end{itemize}
\item {\bf where}: Condiciones que deberá de cumplir la consulta, en lenguaje SQL (lo que habría en la consulta directa a la base de datos tras el ``where''). Se pueden incluir parámetros a introducir previamente a la consulta mediante la sintaxis definida más adelante. Se pueden utilizar las mismas funciones agregadas de base de datos que en el campo show.
\item {\bf order\_by}: Ordenación de los resultados escrito en lenguaje SQL (valor de ``order by'' en la consulta real).
\item {\bf refresh}: Valor de refresco en el caso de consultas en tiempo real. Se puede introducir como parámetro.
\end{itemize}

Los parámetros para Format son los siguientes:
\begin{itemize}
\item {\bf title}: Título del informe. Puede incluir parámetros, por ejemplo, fechas.
\item {\bf rows\_per\_page}: Número de registros a mostrar en cada página en el caso de listas. Puede también definirse como parámetro con valor por omisión, de forma que el usuario puede modificar la opción al solicitar el informe.
\item {\bf width / height}: Ancho y Alto en el caso de gráficas, por ejemplo: ``width=650; height=350;''.
\end{itemize}

\subsection{Parámetros de consulta}
Astatistics permite definir parámetros variables en las plantillas cuyos valores serán solicitados al usuario cuando solicite la visualización de un informe. Pueden ser usados en general en ``Conditions'' y ``Format''. Dichos parámetros pueden ser de los siguientes tipos y se introducen con la siguiente sintaxis:
\begin{itemize}
\item texto: Campo input de formulario web, en el que se puede introducir un texto cualquiera. Su sintaxis es:\\
	\\
	\#text\#nombre a mostrar\#valor por omisión\#\\
	\\
	Para un valor por omisión en blanco se mostrará vacío el campo.\\
	\\
	Un ejemplo podría ser: \#text\#Usuario\#Javier\#\\
	\\
	En algunos casos como el parámetro ``rows\_per\_page'' o ``refresh'' no se ha de especificar el tipo.\\
	\\
	Un ejemplo para ``refresh'' podría ser: \#Refresco\#5\#
\item date: Campo input de formulario que mostrará un desplegable con un calendario para introducción de fechas. Su formato es:\\
	\\
	\#date\#Nombre a mostrar\#valor por omisión\# \\
	\\
	Si no se especifica valor por omisión (\#date\#Nombre\#) se usará de fecha del momento en el que se consulta (fecha de hoy).
\item checkbox: Campo que permite seleccionar uno o varios valores de una lista. Se mostrarán casillas de selección para cada uno de los valores y el usuario podrá marcar los que desee. Su sintaxis es:\\
	\\
\begin{lstlisting}
#checkbox#Nombre(valor1[--checked][|valor2[--checked]...])#
\end{lstlisting}
	
\begin{lstlisting}
Ejemplo:
#checkbox#Usuario(Usuario1--checked,Usuario2,Usuario3--checked)#
\end{lstlisting}
\item radio: Campo similar al anterior pero que solo permite una selección simultánea. La sintaxis es:\\
	\\
\begin{lstlisting}
#radio#Nombre(valor1[--checked][|valor2[--checked]...])#
\end{lstlisting}
	
	Tan solo un valor puede estar marcado como ``checked'' (seleccionado).\\
	
\begin{lstlisting}
Ejemplo:
#radio#Telefono(Tel1,Tel2--checked,Tel3,Tel4)#
\end{lstlisting}
\end{itemize}

Un parámetro puede repetirse para ser utilizado en diferentes partes del diseño de la plantilla. Por ejemplo, un rango de fechas podría ser especificado con \#date\#Fecha inicio\# y \#date\#Fecha fin\# en ``title'' y ``where'' de la siguiente forma:

En Conditions:
\begin{lstlisting}
where=calldate > #date#Fecha inicio# and calldate < #date#Fecha fin#;
\end{lstlisting}
En Format:
\begin{lstlisting}
title=Llamadas entre #date#Fecha inicio# y #date#Fecha fin#;
\end{lstlisting}

Tan solo se solicitaría una vez que el usuario rellenara dichos parámetros al hacer la consulta, usándose luego dicho valor en ambas apariciones en el diseño, es decir, se añadirían dichas fechas al título y se usarían para las condiciones de consulta a la base de datos.

\subsection{Organización de menús e informes}
En ``Templates'' un usuario con rol ``superuser'' o un diseñador podrá organizar plantillas y grupos (otros menús), moverlos, copiarlos, modificarlos y definir cómo se presentarán al usuario final.

Para ello, en ``Templates'' se muestran los grupos y plantillas en estructura jerárquica según la visualización actual, y se muestran botones de operaciones a la derecha de los mismos.

Dicho sistema jerárquico funciona de forma similar a la organización de ficheros en un dispositivo de almacenamiento, como un disco duro o una memoria USB, donde los grupos sería directorios y las plantillas serían ficheros. Dependiendo del menú de trabajo (directorio) encontraremos otros grupos (directorios) o plantillas (ficheros), pudiendo acceder a cualquier grupo visible (directorio) o al grupo padre (bajar un nivel). La ruta se muestra bajo la línea horizontal, justo encima de las plantillas del grupo de trabajo (actual). Es similar a la ruta (path) de un directorio UNIX o Linux, o una página web. Comienza con ``/'' y sería algo similar a ``/grupo1/grupo2/grupo3''.

Al usar el botón de flecha hacia la derecha situado tras un grupo accederemos a su contenido, y el nuevo grupo de trabajo se mostrará en la ruta. Para subir un nivel (volver al grupo padre) hay que usar el botón de flecha hacia la izquierda situado delante de la palabra ``parent'' (padre).

Para mover o duplicar grupos y plantillas en diferentes ubicaciones se usarán los botones ``Copy'' y ``Cut''. Se trata de un sistema ``Cut, copy, paste'' (Cortar, copiar, pegar), clipboard o portapapeles, similar a los utilizados en los escritorios convencionales.

El portapapeles muestra un diálogo al principio de cada página visitada cuando contiene datos. Incluye botones para eliminar el contenido del mismo.

Para cambiar la posición en la lista dentro de un grupo se usarán los botones con flechas verticales a la derecha de grupos y plantillas. La flecha arriba mueve la plantilla o grupo a la posición superior a la actual, mientras que la flecha abajo la mueve hacia abajo.

\chapter{Astatistics: Visualización de informes}
Esta parte del manual de usuario está destinada a los usuarios que deseen visualizar o extraer informes y estadísticas previamente diseñadas por los diseñadores.

Tras acceder a Astatistics (válido para cualquier tipo de usuario) se presenta una lista de las plantillas y grupos disponibles en el grupo raíz ('/').

Tras los grupos se muestra un botón con una flecha hacia la derecha. Dicho botón sirve para acceder al contenido del grupo (entrar).

La flecha hacia la izquierda tras la palabra ``parent'' sirve para acceder al grupo padre (anterior).

Tras las plantillas se muestra un botón ``Execute'', que se deberá de usar para mostrar resultados estadísticos de dicha plantilla.

Tras hacer click en ``Execute'' se mostrará una pantalla intermedia de introducción de datos variables si los hubiera. En dicha pantalla el botón ``parent'' vuelve al menú anterior, y el botón ``Go'' lanza la consulta.

La pantalla de resultados es distinta según el tipo de estadística.

El botón ``Exit'' saldrá de la página de resultados y volverá a la lista anterior a la petición de parámetros.

El botón ``Refresh'' voverá a lanzar la consulta y mostrará los resultados según los datos actuales (que podrían haber variado desde la ejecución anterior).

El botón ``Repeat'' vuelve a solicitar los parámetros de nuevo para la misma consulta.

El botón ``PDF Export'' hace una exportación a formato PDF del resultado actual. En el caso de listas se mostrarán todos los resultados paginados por lo que será necesario repetir la consulta y, dependiendo de lo extensos que sean dichos resultados, puede consumir más tiempo y requerir más recursos generar los resultados (cuando se consulta tan solo una página la consulta especifica el subconjunto de resultados y se efectúa más rápido la consulta a la base de datos).

El botón ``CSV Export'' exportará los resultados en formato CSV (implementado solo para visualización en formato lista en la versión de Astatistics vigente en el momento de escribir estas líneas, aunque se planea posible exportación a CSV de resultados fuente de una gráfica en el futuro. Si se desea dicha opción en versiones que no dispongan de esta función, tan solo es necesario hacer la misma consulta en formato lista y hacer la exportación desde ahí).

\end{document}
